% BOOK-FILE-ID: f-9f1867a090b7
% section: 極形式と複素平面

\section{複素平面と極形式}\label{節：複素平面と極形式}
複素数の積を回転と拡大・縮小として読む.
\sectionref*{節：複素平面で見る}と\sectionref*{節：特性方程式の解が虚数の場合}で使う.
\bookterm{polar}{極形式}には\bookterm{trigonometric}{三角関数},\bookterm{euler}{オイラーの公式}の導出には\bookterm{limit}{極限}を使う.

\subsection{複素数を平面に表す}\label{節：複素数を平面に表す}
$i^2=-1$を満たす虚数単位$i$を用い,複素数$z=x+iy$を点$(x,y)$に対応させたものを\bookterm{complex-plane}{複素平面}という.
横軸が実軸,縦軸が虚軸である.
$z=x+iy$に対して,\,$\overline z=x-iy$を$z$の共役複素数という.
複素平面上では実軸について対称な点であり,\,$|\overline z|=|z|$である.
また,\,$\overline{z+w}=\overline z+\overline w$,\,$\overline{zw}=\overline z\,\overline w$が成り立つ.
したがって,係数が実数の式では,\,$z$を$\overline z$に変えると式の値も共役になる.
原点からの距離を$\rho=\sqrt{x^2+y^2}$,正の実軸から反時計回りに測る角を$\theta$とすると,
\[
  x=\rho\cos\theta,\qquad y=\rho\sin\theta
\]
より
\[
  z=\rho(\cos\theta+i\sin\theta).
\]
これが\bookterm{polar}{極形式}である.\,$\rho=|z|$は絶対値,\,$z\ne0$のとき$\theta$は\bookterm{argument}{偏角}である.
偏角に$2\pi$の整数倍を加えても同じ点を表す.\,$z=0$では偏角は定まらない.
\BookFigPolar

例えば$\rho=2$, $\theta=\frac{\pi}{3}$なら
\[
  z=2\left(\cos\frac\pi3+i\sin\frac\pi3\right)=1+i\sqrt3
\]
であり,点$(1,\sqrt3)$を表す.

\subsection{複素数と平面の変換}\label{節：複素数と平面の変換}
$b=u+iv$を加えると,\,$z+b=(x+u)+i(y+v)$となる.
これはベクトル$(u,v)$だけの平行移動である.
次に,\,$c=s(\cos\phi+i\sin\phi)$, $s>0$を掛ける.
\bookterm{addition}{加法定理}より
\[
  cz=s\rho\{\cos(\theta+\phi)+i\sin(\theta+\phi)\}.
\]
絶対値は$s$倍,偏角は$\phi$だけ増える.特に$|cz|=|c||z|$であり,\,$w\ne0$なら$\left|\frac zw\right|=\frac{|z|}{|w|}$である.原点を中心に回転し,\,$s>1$なら拡大,\,$0<s<1$なら縮小,\,$s=1$なら回転だけになる.
\BookFigComplexRotation

特に$i$を掛けると反時計回りに$\frac{\pi}{2}$回転する.
中心を$z_0$から$w_0$へ移す変換は
\[
  w=c(z-z_0)+w_0.
\]
中心からの差に$c$を掛け,最後に移動先の中心を加える.

\subsection{ド・モアブルの定理}\label{節：ド・モアブルの定理}
先の積の公式では,絶対値が$1$なら偏角だけが加わる.繰り返すと
\[
  (\cos\theta+i\sin\theta)^n=\cos(n\theta)+i\sin(n\theta)
\]
を得る.これを\bookterm{de-moivre}{ド・モアブルの定理}という.
正の整数$n$では\bookterm{induction}{数学的帰納法}で示せる.
$n=1$で一致し,\,$n$での式に$\cos\theta+i\sin\theta$を掛けると,加法定理で角度が$n\theta$から$(n+1)\theta$になる.
$n=0$では両辺とも$1$である.
また,
\[
  (\cos\theta+i\sin\theta)^{-1}=\cos(-\theta)+i\sin(-\theta)
\]
なので,負の整数も正の整数の場合へ帰着できる.
例えば$\theta=\frac{\pi}{6}$, $n=6$なら$\cos\pi+i\sin\pi=-1$となる.

\subsection{オイラーの公式}\label{節：オイラーの公式}
複素数の指数表示を
\[
  e^{i\theta}=\lim_{n\to\infty}\left(1+\frac{i\theta}{n}\right)^n
\]
と定め,\,\sectionref{節：極限}の基本極限から値を求める.
$\theta\ne0$とし,\,$-\frac{\pi}{2}<\phi_n<\frac{\pi}{2}$, $\tan\phi_n=\frac{\theta}{n}$とおく.
$r_n=\sqrt{1+\dfrac{\theta^2}{n^2}}$なら
\[
  1+\frac{i\theta}{n}=r_n(\cos\phi_n+i\sin\phi_n).
\]
長さと角度の極限を別々に求める.
\begin{align*}
  \log(r_n^n)
  &=\frac{\theta^2}{2n}\,
  \cfrac{\log\left(1+\cfrac{\theta^2}{n^2}\right)}{\cfrac{\theta^2}{n^2}}\longrightarrow0,\\
  n\phi_n&=\theta\frac{\phi_n}{\tan\phi_n}\longrightarrow\theta.
\end{align*}
最初の式では$\frac{\log(1+x)}{x}\to1$,後では$\phi_n\to0$と$\frac{x}{\tan x}\to1$を使った.
$r_n^n\to1$とド・モアブルの定理から
\[
  \left(1+\frac{i\theta}{n}\right)^n
  =r_n^n\{\cos(n\phi_n)+i\sin(n\phi_n)\}.
\]
三角関数の連続性により
\[
  e^{i\theta}=\cos\theta+i\sin\theta
\]
を得る.\,$\theta=0$でも両辺は$1$である.これが\bookterm{euler}{オイラーの公式}である.

極形式は$z=\rho e^{i\theta}$と短く書ける.\,$z\ne0$,すなわち$\rho>0$なら,逆数を含む整数乗も
\[
  \begin{aligned}
    z^n&=\rho^n e^{in\theta}\\
    &=\rho^n\{\cos(n\theta)+i\sin(n\theta)\}
    \quad(n\in\mathbb Z)
  \end{aligned}
\]
と表せる.\,$z=0$では負の整数乗を定めず,正の整数$n$について$0^n=0$とする.

\subsection{複素数の実数乗}\label{節：複素数の実数乗}
$z=\rho e^{i\theta}$, $\rho>0$とし,偏角$\theta$を一つ固定すれば,実数$x$について
\[
  z^x=\rho^x e^{ix\theta}
\]
と定められる.ただし,\,$\theta$を$\theta+2\pi k$へ変えると,整数でない$x$では値が変わりうる.
整数なら偏角の選び方によらない.
有理数$x=\frac{m}{n}$, $n\ge1$では$\left(z^{\dfrac{m}{n}}\right)^n=z^m$となり,\,$z^{\dfrac{m}{n}}$は$z^m$の$n$乗根の一つである.

\subsection{公式のまとめ}\label{節：公式のまとめ}
\bookterm{polar}{極形式}$z=\rho e^{i\theta}$では,加法は平行移動,乗法は長さの積と角度の和,整数乗は長さの累乗と角度の整数倍に対応する.
複素数の計算を平面の変換として読めるのである.
